% Pursuing Stacks (cotes 134-2 to 134-8): archive cote 134-3, site page 12 (PDF page 13), typescript page 12.
% typescript. Draft compared: ps2.tex:516-584.
% AI-transcribed by Claude Opus 5.5 (claude-opus-5-5, adaptive thinking, effort high), Message Batches API, batch ps1 (msgbatch_01Rv5hCaNTtUt4JHx1yqoHPL), 2026-11-12; not yet checked by eye.
\scanpage{134-3/12}
\mspage{12}

\hand{19.} I would like now to write down a provisional itinerary of the \unsure{voyage}
ahead - namely to make a list of those main features of a theory of
\margin{$\times$} stacks which are in my mind these days. I will write them down in the
order in which they occur to me - which will be no obligation upon me
to follow this order, when coming back on those features separately to
elaborate somewhat on them\over{,}{.} \inserted{This I expect to do,} mainly as a way to check wether the main
notions and intuitions introduced are sound indeed, and otherwise, to
see how to correct them.

1\textsuperscript{o}) Definition of the categories $\underline{B}_o$ and its dual $\underline{C}_o$ , and formal defi-%
nition of left and right coherators. This definition \over{\unsure{will}}{will} still be a
provisional one, and will presumably have to be adjusted somewhat to
\over{allo}{allo}w for the various structures we are looking for in the corresponding
category of Gr-stacks.

2\textsuperscript{o}) Relation between the category of Gr-stacks and the category of
topological spaces, via two adjoint functors, the ``topological realiza-%
tion functor'' $F_{\ast} \hand{$\mapsto$} \hand{$|$}F_{\ast}\hand{$|$}$ , and the ``singular stack functor''
$X \hand{$\mapsto$} \underline{F}_{\ast}(X)$ . The situation should be formally analogous to the corres-%
ponding situation for semi-simplicial sets versus topological spaces,
the role of the category $\underline{S}_{\ast}$ of standar\over{t}{d} ordered simplices being taken
by the left coherator $\underline{B}$ we are working with. The main technical diffe-%
rence here is that the category of Gr-stacks is not just the category of
presheaves on $\underline{B}$, but the full subcategory defined by the requirement
that the presheaves considered should transform ``standar\over{t}{d}'' amalgamated
sums into fibered products. As a matter of fact, the topological rea-%
lization functor $F_{\ast} \hand{$\mapsto$} \hand{$|$}F_{\ast}\hand{$|$}$ could be defined in standar\over{t}{d} way on the
whole of $\hand{$\widehat{\phantom{B}}$}\underline{B}$ , in terms of any functor \xout{B}
\hand{(*)} \[ \underline{B} \hand{$\longrightarrow$} \text{(Spaces)} \]
(by the requirement that the extension of this functor to $\hand{$\widehat{\phantom{B}}$}B$ commutes
with arbitrary direct limits). A second difference with the \inserted{similar,} simplicial
\deleted{similar} situation \xout{that} lies in the fact that the only really compel-%
ling choice for the functor (*) is it's restriction to $\underline{B}_o$ , in terms
of the cells $\underline{D}_n$ and the standart \hand{``}half-hemisphere \xout{morphisms} maps\hand{''} between
these (L, \deleted{\illegible{1}}\inserted{section 6}). The extension of this functor to $\underline{B}$ is always possi\hand{\unsure{b}}%
le, due to the \inserted[typed]{inductive} construction of $\underline{B}$ and to the interpretation of elements
\margin{(*)$_o$ (circled)} of $\underline{B}_o$ as \uline{contractible} spaces, via the functor (*)\hand{$_o$}; but it depends on
a bunch of arbitrary choices. To give precise meaning to the intuition
that these choices do\over{nt}{'nt} really make a difference, and that the choic\unsure{e}
\inserted{of} coherator we are starting with does'nt make much of a difference either\hand{,}
will need some elaboration on the notion of \uline{equivalences between}
